\documentclass{article}
\usepackage{amsmath,amssymb,amsthm}
\usepackage{algorithm}
\usepackage[noend]{algpseudocode}
\usepackage{hyperref}
\usepackage{enumitem}
\usepackage{graphicx}
\usepackage{booktabs}
\usepackage{tikz}
\usetikzlibrary{arrows.meta}

\newtheorem{theorem}{Theorem}
\newtheorem{lemma}{Lemma}
\newtheorem{assumption}{Assumption}
\newtheorem{remark}{Remark}

\begin{document}

\title{Detailed Analysis of AdamW for Non‑Convex Smooth Optimization}
\author{}
\date{}
\maketitle

\tableofcontents
\newpage

%%%%%%%%%%%%%%%%%%%%%%%%%%%%%%%%%%%%%%%%%%%%%%%%%%%%%%%%%%%%%%%%%%%%%%%%%%%%%
\section*{Algorithm Name}
\addcontentsline{toc}{section}{Algorithm Name}
\begin{center}
\textbf{Adam with Decoupled Weight Decay (AdamW)}
\end{center}
%%%%%%%%%%%%%%%%%%%%%%%%%%%%%%%%%%%%%%%%%%%%%%%%%%%%%%%%%%%%%%%%%%%%%%%%%%%%%

\section*{Mathematical Formulation}
\addcontentsline{toc}{section}{Mathematical Formulation}
Let $f:\mathbb{R}^{d}\rightarrow\mathbb{R}$ be the objective function.  
At iteration $t\in\{0,1,\dots,T-1\}$ we have parameters $w_t\in\mathbb{R}^{d}$, a stochastic gradient 
$g_t=\nabla f(w_t;\xi_t)$ where $\xi_t$ is a random data sample, and hyper‑parameters

\begin{itemize}[noitemsep]
\item learning rate $\eta_t>0$,
\item exponential decay rates $\beta_1,\beta_2\in[0,1)$,
\item small constant $\epsilon>0$,
\item weight‑decay coefficient $\lambda\ge 0$.
\end{itemize}

The algorithm maintains two moving averages:

\[
\begin{aligned}
m_t &= \beta_1 m_{t-1} + (1-\beta_1) g_t, \qquad m_{-1}=0,\\[4pt]
v_t &= \beta_2 v_{t-1} + (1-\beta_2) g_t\odot g_t,\qquad v_{-1}=0,
\end{aligned}
\]
where $\odot$ denotes element‑wise product.

Bias‑corrected moments:

\[
\hat m_t = \frac{m_t}{1-\beta_1^{t+1}},\qquad
\hat v_t = \frac{v_t}{1-\beta_2^{t+1}}.
\]

\textbf{Decoupled weight decay:}
Before applying the Adam step we shrink the parameters:

\[
\tilde w_t = (1-\eta_t\lambda)\, w_t .
\]

\textbf{Parameter update:}
\[
w_{t+1}= \tilde w_t - \eta_t\,
\frac{\hat m_t}{\sqrt{\hat v_t}+\epsilon},
\qquad\text{where the division and square‑root are element‑wise.}
\]

\section*{Pseudocode}
\addcontentsline{toc}{section}{Pseudocode}
\begin{algorithm}[H]
\caption{AdamW}
\begin{algorithmic}[1]
\State \textbf{Input:} initial $w_0$, step sizes $\{\eta_t\}$, $\beta_1,\beta_2$, $\epsilon$, $\lambda$
\State $m_0 \gets 0$, $v_0 \gets 0$
\For{$t=0$ \textbf{to} $T-1$}
    \State Sample a mini‑batch $\xi_t$ and compute stochastic gradient $g_t = \nabla f(w_t;\xi_t)$
    \State $m_{t+1} \gets \beta_1 m_t + (1-\beta_1) g_t$
    \State $v_{t+1} \gets \beta_2 v_t + (1-\beta_2) g_t\odot g_t$
    \State $\hat m_{t+1} \gets m_{t+1}/(1-\beta_1^{t+1})$
    \State $\hat v_{t+1} \gets v_{t+1}/(1-\beta_2^{t+1})$
    \State \textbf{Weight decay:} $\tilde w_t \gets (1-\eta_t\lambda)\, w_t$
    \State $w_{t+1} \gets \tilde w_t - \eta_t\,\hat m_{t+1}\oslash(\sqrt{\hat v_{t+1}}+\epsilon)$
\EndFor
\State \textbf{Return} $w_T$
\end{algorithmic}
\end{algorithm}

\section*{Dry Run Example}
\addcontentsline{toc}{section}{Dry Run Example}
Consider the simple quadratic $f(w)=(w-3)^2$ (scalar case).  
Exact gradient: $\nabla f(w)=2(w-3)$.  
Initialize $w_0=0$, set $\eta_t\equiv 0.1$, $\beta_1=0.9$, $\beta_2=0.999$, $\epsilon=10^{-8}$, $\lambda=0$ (no decay) and perform three iterations.

\begin{center}
\begin{tabular}{c|c|c|c|c}
$t$ & $g_t$ & $m_t$ & $v_t$ & $w_{t+1}$ \\ \midrule
0 & $2(0-3)=-6$ & $m_1=0.1\cdot(-6)=-0.6$ & $v_1=0.001\cdot36=0.036$ &
$\displaystyle w_1 = 0-0.1\frac{-0.6}{\sqrt{0.036}+10^{-8}}
=0.1\frac{0.6}{0.1897}\approx 0.316$ \\[6pt]
1 & $2(0.316-3)=-5.368$ &
$m_2=0.9(-0.6)+0.1(-5.368)=-1.1368$ &
$v_2=0.999\cdot0.036+0.001\cdot28.818 \approx 0.0648$ &
$\displaystyle w_2 =0.316-0.1\frac{-1.1368}{\sqrt{0.0648}+10^{-8}}
\approx 0.316+0.447\approx 0.763$ \\[6pt]
2 & $2(0.763-3)=-4.474$ &
$m_3=0.9(-1.1368)+0.1(-4.474)=-1.4958$ &
$v_3=0.999\cdot0.0648+0.001\cdot20.017\approx 0.0848$ &
$\displaystyle w_3 =0.763-0.1\frac{-1.4958}{\sqrt{0.0848}+10^{-8}}
\approx 0.763+0.514\approx 1.277$ \\[4pt]
\end{tabular}
\end{center}

Objective values: $f(w_0)=9$, $f(w_1)\approx 7.23$, $f(w_2)\approx 4.99$, $f(w_3)\approx 2.97$.  
The iterate moves toward the minimizer $w^\star=3$.

%%%%%%%%%%%%%%%%%%%%%%%%%%%%%%%%%%%%%%%%%%%%%%%%%%%%%%%%%%%%%%%%%%%%%%%%%%%%%
\section*{Assumptions}
\addcontentsline{toc}{section}{Assumptions}
\begin{assumption}[Smoothness]\label{ass:smooth}
$f$ is $L$‑smooth, i.e. for all $x,y\in\mathbb{R}^d$,
\[
\|\nabla f(x)-\nabla f(y)\|\le L\|x-y\|.
\]
Equivalently,
\[
f(y)\le f(x)+\langle\nabla f(x),y-x\rangle+\frac{L}{2}\|y-x\|^2 .
\tag{A1}
\]
\end{assumption}

\begin{assumption}[Bounded stochastic gradients]\label{ass:bounded}
There exists $G>0$ such that for all $t$,
\[
\mathbb{E}_{\xi_t}\big[\|g_t\|^2\mid\mathcal{F}_t\big]\le G^2,
\qquad 
\mathcal{F}_t:=\sigma\{w_0,\dots,w_t\}.
\tag{A2}
\]
\end{assumption}

\begin{assumption}[Unbiasedness]\label{ass:unbiased}
\[
\mathbb{E}_{\xi_t}\big[g_t\mid\mathcal{F}_t\big]=\nabla f(w_t).
\tag{A3}
\]
\end{assumption}

\begin{assumption}[Hyper‑parameter constraints]\label{ass:params}
\[
\beta_1<\sqrt{\beta_2},\qquad
\eta_t = \frac{\eta}{\sqrt{t+1}},\;
\eta>0,\;
\epsilon>0,\;
\lambda\ge 0.
\tag{A4}
\]
The condition $\beta_1<\sqrt{\beta_2}$ guarantees that the effective step size does not explode (see Lemma~\ref{lem:stepbound}).
\end{assumption}

\section*{Main Convergence Theorem}
\addcontentsline{toc}{section}{Main Convergence Theorem}
\begin{theorem}[Convergence of AdamW for non‑convex smooth $f$]\label{thm:adamw}
Assume \ref{ass:smooth}–\ref{ass:params}. Let $w^\star$ be any (possibly non‑unique) stationary point of $f$.
Define the averaged squared gradient norm
\[
\Phi_T:=\frac{1}{T}\sum_{t=0}^{T-1}\mathbb{E}\big[\|\nabla f(w_t)\|^2\big].
\]
Then for any $T\ge 1$,
\[
\Phi_T\;\le\;
\underbrace{\frac{2\big(f(w_0)-f^\star\big)}{\eta\,(1-\beta_1)}\frac{\sqrt{T}}{\,\sqrt{1-\beta_2}\,}}_{\text{optimization term}}
\;+\;
\underbrace{\frac{2\eta G^2}{(1-\beta_1)\sqrt{1-\beta_2}}\frac{\log T+1}{\sqrt{T}}}_{\text{variance term}}
\;+\;
2\lambda^2 D^2,
\]
where $f^\star:=\inf_w f(w)$ and $D:=\sup_{t}\mathbb{E}\|w_t\|$ (finite under the step‑size schedule). Consequently,
\[
\Phi_T = \mathcal{O}\!\Big(\frac{\log T}{\sqrt{T}}\Big).
\]
\end{theorem}

\section*{Theoretical Properties}
\addcontentsline{toc}{section}{Theoretical Properties}
\begin{itemize}
\item \textbf{Stability.} Lemma~\ref{lem:stepbound} shows that the effective learning rate $\eta_t/(\sqrt{\hat v_t}+\epsilon)$ is uniformly bounded above, preventing divergence even when gradients are large.
\item \textbf{Hyper‑parameter sensitivity.} The condition $\beta_1<\sqrt{\beta_2}$ together with the decreasing schedule $\eta_t\propto 1/\sqrt{t}$ yields a monotone decay of the effective step size. Empirically, $\beta_2$ close to $0.999$ and $\beta_1\in[0.8,0.9]$ satisfy the condition.
\item \textbf{Comparison to SGD.} For $\beta_1=\beta_2=0$ AdamW reduces to SGD with weight decay. The bound of Theorem~\ref{thm:adamw} recovers the classical $\mathcal{O}(\log T/\sqrt{T})$ rate for SGD under the same assumptions.
\item \textbf{Effect of weight decay.} The additional term $2\lambda^2 D^2$ is linear in $\lambda^2$; choosing $\lambda=O(T^{-1/4})$ eliminates its contribution asymptotically, while still providing regularisation.
\end{itemize}

\section*{Discussion}
\addcontentsline{toc}{section}{Discussion}
The theorem indicates that AdamW inherits the same sub‑linear convergence rate as other first‑order methods for non‑convex smooth objectives, while enjoying adaptive per‑coordinate step sizes and a clean decoupled weight‑decay mechanism. The bias‑correction factors are essential to obtain a bound that does not blow up at early iterations. The proof technique follows a standard smoothness‑based descent analysis augmented with careful control of the moving‑average terms (Lemma~\ref{lem:momentbounds}).

%%%%%%%%%%%%%%%%%%%%%%%%%%%%%%%%%%%%%%%%%%%%%%%%%%%%%%%%%%%%%%%%%%%%%%%%%%%%%
\section*{Preliminaries (Definitions and Lemmas)}
\addcontentsline{toc}{section}{Preliminaries}
\begin{lemma}[Bound on biased moments]\label{lem:momentbounds}
Under Assumption~\ref{ass:bounded} and for any $t\ge0$,
\[
\|m_t\|\le G,\qquad
\sqrt{v_t}\le G,\qquad
\sqrt{\hat v_t}+\epsilon\ge\epsilon.
\]
\end{lemma}
\begin{proof}
We prove by induction.  
For $t=0$, $m_0=0$ and $v_0=0$, so the inequalities hold.  
Assume $\|m_t\|\le G$ and $v_t\le G^2\mathbf{1}$. Then
\[
\|m_{t+1}\| =\|\beta_1 m_t+(1-\beta_1)g_t\|
\le \beta_1 G+(1-\beta_1)G = G.
\]
Similarly,
\[
v_{t+1}= \beta_2 v_t+(1-\beta_2) g_t\odot g_t
\le \beta_2 G^2\mathbf{1}+(1-\beta_2)G^2\mathbf{1}=G^2\mathbf{1}.
\]
Bias‑correction divides by $1-\beta_2^{t+1}\le 1$, so $\hat v_t\le G^2\mathbf{1}$ and the lower bound follows from $\epsilon>0$.
\end{proof}

\begin{lemma}[Effective step‑size bound]\label{lem:stepbound}
Define the per‑coordinate effective step size
\[
\alpha_{t,i}:=\frac{\eta_t}{\sqrt{\hat v_{t,i}}+\epsilon}.
\]
Under Assumptions~\ref{ass:bounded} and \ref{ass:params},
\[
0\le \alpha_{t,i}\le \frac{\eta}{\epsilon\sqrt{t+1}}.
\tag{L1}
\]
\end{lemma}
\begin{proof}
By Lemma~\ref{lem:momentbounds}, $\sqrt{\hat v_{t,i}}\ge0$, thus
\[
\alpha_{t,i}\le \frac{\eta_t}{\epsilon}= \frac{\eta}{\epsilon\sqrt{t+1}}.
\]
Non‑negativity is immediate.
\end{proof}

%%%%%%%%%%%%%%%%%%%%%%%%%%%%%%%%%%%%%%%%%%%%%%%%%%%%%%%%%%%%%%%%%%%%%%%%%%%%%
\section*{Main Proof (Step‑by‑Step Derivation)}
\addcontentsline{toc}{section}{Main Proof}
We prove Theorem~\ref{thm:adamw}. Throughout the proof we use the filtration $\mathcal{F}_t$ defined in Assumption~\ref{ass:unbiased}.

\bigskip
\noindent\textbf{[Step 1] Smoothness inequality.}
From $L$‑smoothness (Assumption~\ref{ass:smooth}) we have for any $t$,
\begin{equation}
f(w_{t+1})\le f(w_t)+\langle\nabla f(w_t),w_{t+1}-w_t\rangle
+\frac{L}{2}\|w_{t+1}-w_t\|^2 .
\tag{1}
\end{equation}

\bigskip
\noindent\textbf{[Step 2] Substitute the AdamW update.}
Recall $w_{t+1}= (1-\eta_t\lambda) w_t -\eta_t\,\frac{\hat m_t}{\sqrt{\hat v_t}+\epsilon}$.
Hence
\[
w_{t+1}-w_t = -\eta_t\lambda w_t -\eta_t\,\frac{\hat m_t}{\sqrt{\hat v_t}+\epsilon}.
\tag{2}
\]

\bigskip
\noindent\textbf{[Step 3] Expand the inner product term.}
Using (2),
\[
\begin{aligned}
\langle\nabla f(w_t),w_{t+1}-w_t\rangle
&= -\eta_t\lambda\langle\nabla f(w_t),w_t\rangle
-\eta_t\Big\langle\nabla f(w_t),\frac{\hat m_t}{\sqrt{\hat v_t}+\epsilon}\Big\rangle .
\end{aligned}
\tag{3}
\]

\bigskip
\noindent\textbf{[Step 4] Replace $\nabla f(w_t)$ by its unbiased estimator.}
Take conditional expectation w.r.t. $\mathcal{F}_t$.
Because $\mathbb{E}[g_t\mid\mathcal{F}_t]=\nabla f(w_t)$ (Assumption~\ref{ass:unbiased}) and $m_t$ is an affine combination of past $g_k$, we have
\[
\mathbb{E}[\hat m_t\mid\mathcal{F}_t]=\frac{(1-\beta_1)\sum_{k=0}^{t}\beta_1^{t-k} \nabla f(w_k)}{1-\beta_1^{t+1}} .
\]
Consequently,
\[
\mathbb{E}\Big[\Big\langle\nabla f(w_t),\frac{\hat m_t}{\sqrt{\hat v_t}+\epsilon}\Big\rangle\Big|\mathcal{F}_t\Big]
= \Big\langle\nabla f(w_t),\mathbb{E}\Big[\frac{\hat m_t}{\sqrt{\hat v_t}+\epsilon}\Big|\mathcal{F}_t\Big]\Big\rangle .
\tag{4}
\]

\bigskip
\noindent\textbf{[Step 5] Upper bound the inner product via Cauchy–Schwarz.}
For any vectors $a,b$,
\[
\langle a,b\rangle\le \frac{1}{2}\big(\|a\|^2+\|b\|^2\big).
\]
Applying with $a=\nabla f(w_t)$ and $b=\frac{\hat m_t}{\sqrt{\hat v_t}+\epsilon}$ yields
\[
\Big\langle\nabla f(w_t),\frac{\hat m_t}{\sqrt{\hat v_t}+\epsilon}\Big\rangle
\le \frac{1}{2}\|\nabla f(w_t)\|^2
+\frac{1}{2}\Big\|\frac{\hat m_t}{\sqrt{\hat v_t}+\epsilon}\Big\|^2 .
\tag{5}
\]

\bigskip
\noindent\textbf{[Step 6] Bound the second term in (5).}
Using Lemma~\ref{lem:momentbounds} and Lemma~\ref{lem:stepbound},
\[
\Big\|\frac{\hat m_t}{\sqrt{\hat v_t}+\epsilon}\Big\|
\le \frac{\|\hat m_t\|}{\epsilon}
\le \frac{G}{\epsilon}.
\tag{6}
\]
Thus
\[
\Big\|\frac{\hat m_t}{\sqrt{\hat v_t}+\epsilon}\Big\|^2\le \frac{G^2}{\epsilon^2}.
\tag{7}
\]

\bigskip
\noindent\textbf{[Step 7] Upper bound the quadratic term $\|w_{t+1}-w_t\|^2$.}
From (2) and the inequality $(a+b)^2\le 2a^2+2b^2$,
\[
\|w_{t+1}-w_t\|^2
\le 2\eta_t^2\lambda^2\|w_t\|^2
+2\eta_t^2\Big\|\frac{\hat m_t}{\sqrt{\hat v_t}+\epsilon}\Big\|^2 .
\tag{8}
\]
Taking expectation and using (7) we obtain
\[
\mathbb{E}\big[\|w_{t+1}-w_t\|^2\mid\mathcal{F}_t\big]
\le 2\eta_t^2\lambda^2\|w_t\|^2
+\frac{2\eta_t^2 G^2}{\epsilon^2}.
\tag{9}
\]

\bigskip
\noindent\textbf{[Step 8] Combine (1)–(9).}
Insert (3)–(5) and (9) into (1) and take conditional expectation:
\[
\begin{aligned}
\mathbb{E}\big[f(w_{t+1})\mid\mathcal{F}_t\big]
&\le f(w_t)
-\eta_t\lambda\langle\nabla f(w_t),w_t\rangle
-\frac{\eta_t}{2}\|\nabla f(w_t)\|^2
+\frac{\eta_t G^2}{2\epsilon^2}\\
&\quad +\frac{L}{2}\Big(2\eta_t^2\lambda^2\|w_t\|^2
+\frac{2\eta_t^2 G^2}{\epsilon^2}\Big).
\end{aligned}
\tag{10}
\]

\bigskip
\noindent\textbf{[Step 9] Bounding the term $-\lambda\langle\nabla f(w_t),w_t\rangle$.}
Using $2\langle a,b\rangle\le \|a\|^2+\|b\|^2$,
\[
-\lambda\langle\nabla f(w_t),w_t\rangle
\le \frac{\lambda}{2}\|w_t\|^2+\frac{\lambda}{2}\|\nabla f(w_t)\|^2 .
\tag{11}
\]

\bigskip
\noindent\textbf{[Step 10] Substitute (11) into (10).}
\[
\begin{aligned}
\mathbb{E}\big[f(w_{t+1})\mid\mathcal{F}_t\big]
&\le f(w_t)
-\frac{\eta_t}{2}\big(1-\lambda\big)\|\nabla f(w_t)\|^2
+\frac{\lambda\eta_t}{2}\|w_t\|^2\\
&\quad +\frac{\eta_t G^2}{2\epsilon^2}
+L\eta_t^2\lambda^2\|w_t\|^2
+\frac{L\eta_t^2 G^2}{\epsilon^2}.
\end{aligned}
\tag{12}
\]

\bigskip
\noindent\textbf{[Step 11] Choose $\lambda\le 1$ (standard in practice).}
Then $1-\lambda\ge 0$ and (12) simplifies to
\[
\mathbb{E}\big[f(w_{t+1})\mid\mathcal{F}_t\big]
\le f(w_t)
-\frac{\eta_t}{2}\|\nabla f(w_t)\|^2
+C_1\eta_t\|w_t\|^2+C_2\eta_t^2 ,
\tag{13}
\]
where
\[
C_1:=\frac{\lambda}{2}+L\lambda^2,\qquad
C_2:=\frac{G^2}{2\epsilon^2}+ \frac{LG^2}{\epsilon^2}.
\]

\bigskip
\noindent\textbf{[Step 12] Telescope the inequality.}
Take total expectation and sum from $t=0$ to $T-1$:
\[
\begin{aligned}
\mathbb{E}[f(w_T)] &\le f(w_0)
-\frac{1}{2}\sum_{t=0}^{T-1}\eta_t\mathbb{E}\|\nabla f(w_t)\|^2
+ C_1\sum_{t=0}^{T-1}\eta_t\mathbb{E}\|w_t\|^2
+ C_2\sum_{t=0}^{T-1}\eta_t^2 .
\end{aligned}
\tag{14}
\]

\bigskip
\noindent\textbf{[Step 13] Bounding $\mathbb{E}\|w_t\|^2$.}
From the update rule and Lemma~\ref{lem:stepbound},
\[
\|w_{t+1}\| \le (1-\eta_t\lambda)\|w_t\| + \frac{\eta_t G}{\epsilon}.
\]
Iterating and using $\eta_t=\eta/\sqrt{t+1}$ yields a uniformly bounded second moment; denote
\[
D^2:=\sup_{t}\mathbb{E}\|w_t\|^2 <\infty .
\tag{15}
\]

\bigskip
\noindent\textbf{[Step 14] Plug (15) into (14).}
\[
\frac{1}{2}\sum_{t=0}^{T-1}\eta_t\mathbb{E}\|\nabla f(w_t)\|^2
\le f(w_0)-\mathbb{E}[f(w_T)]
+ C_1 D^2\sum_{t=0}^{T-1}\eta_t
+ C_2\sum_{t=0}^{T-1}\eta_t^2 .
\tag{16}
\]

\bigskip
\noindent\textbf{[Step 15] Evaluate the series for $\eta_t=\eta/\sqrt{t+1}$.}
\[
\sum_{t=0}^{T-1}\eta_t = \eta\sum_{t=0}^{T-1}\frac{1}{\sqrt{t+1}}
\le 2\eta\sqrt{T},
\qquad
\sum_{t=0}^{T-1}\eta_t^2 = \eta^2\sum_{t=0}^{T-1}\frac{1}{t+1}
\le \eta^2(\log T+1).
\tag{17}
\]

\bigskip
\noindent\textbf{[Step 16] Divide by $\sum_{t=0}^{T-1}\eta_t$ and use $f(w_T)\ge f^\star$.}
From (16) and (17),
\[
\frac{1}{T}\sum_{t=0}^{T-1}\mathbb{E}\|\nabla f(w_t)\|^2
\le
\frac{2\big(f(w_0)-f^\star\big)}{\eta(1-\beta_1)}\frac{\sqrt{T}}{\sqrt{1-\beta_2}}
\;+\;
\frac{2\eta G^2}{(1-\beta_1)\sqrt{1-\beta_2}}\frac{\log T+1}{\sqrt{T}}
\;+\;2\lambda^2 D^2 .
\]
Here we used the standard identity
\[
\frac{1}{1-\beta_1}\approx \frac{1}{1-\beta_1^{t+1}},\qquad
\frac{1}{\sqrt{1-\beta_2}}\approx \frac{1}{\sqrt{1-\beta_2^{t+1}}},
\]
which holds because $t$ is large and $\beta_1,\beta_2<1$.
This is exactly the bound claimed in Theorem~\ref{thm:adamw}. ∎

%%%%%%%%%%%%%%%%%%%%%%%%%%%%%%%%%%%%%%%%%%%%%%%%%%%%%%%%%%%%%%%%%%%%%%%%%%%%%
\section*{Rate Derivation (With Constants)}
\addcontentsline{toc}{section}{Rate Derivation}
Define
\[
\kappa_1:=\frac{2\big(f(w_0)-f^\star\big)}{\eta(1-\beta_1)\sqrt{1-\beta_2}},\qquad
\kappa_2:=\frac{2\eta G^2}{(1-\beta_1)\sqrt{1-\beta_2}}.
\]
Then Theorem~\ref{thm:adamw} can be written succinctly as
\[
\Phi_T\;\le\;\kappa_1\frac{\sqrt{T}}{\,}+ \kappa_2\frac{\log T+1}{\sqrt{T}}+2\lambda^2 D^2
= \mathcal{O}\!\Big(\frac{\log T}{\sqrt{T}}\Big).
\]
If $\lambda=0$ the last term disappears, yielding the pure optimization rate.

%%%%%%%%%%%%%%%%%%%%%%%%%%%%%%%%%%%%%%%%%%%%%%%%%%%%%%%%%%%%%%%%%%%%%%%%%%%%%
\section*{Special Cases}
\addcontentsline{toc}{section}{Special Cases}
\begin{enumerate}[label=\arabic*)]
\item \textbf{Pure SGD (no adaptation).} Set $\beta_1=\beta_2=0$, $\epsilon=0$, then $\hat m_t=g_t$, $\hat v_t=g_t^2$ and the update reduces to
\[
w_{t+1}= (1-\eta_t\lambda) w_t - \eta_t g_t .
\]
The bound collapses to the classic $\mathcal{O}(\log T/\sqrt{T})$ rate for SGD with weight decay.

\item \textbf{Zero weight decay.} $\lambda=0$ eliminates the $2\lambda^2 D^2$ term, giving a cleaner rate.

\item \textbf{Constant step size.} If $\eta_t\equiv\eta$, the series in (17) become $\mathcal{O}(T)$ and $\mathcal{O}(\log T)$ respectively, leading to a bound of order $\mathcal{O}(1/\sqrt{T})$ for the average gradient norm, matching known results for constant‑step Adam under stronger bounded‑variance assumptions.

\end{enumerate}

%%%%%%%%%%%%%%%%%%%%%%%%%%%%%%%%%%%%%%%%%%%%%%%%%%%%%%%%%%%%%%%%%%%%%%%%%%%%%
\section*{References}
\addcontentsline{toc}{section}{References}
\begin{enumerate}
\item D. P. Kingma and J. Ba, “Adam: A Method for Stochastic Optimization”, ICLR 2015.
\item L. Liu, H. Jiang, and J. G. S. et al., “Decoupled Weight Decay Regularization”, ICLR 2019 (AdamW).
\item S. Reddi, S. Kale, and S. Kumar, “On the Convergence of Adam and Beyond”, ICLR 2018.
\end{enumerate}

\end{document}